\documentclass[11pt,a4paper]{article}
\usepackage[utf8]{inputenc}
\usepackage[ngerman]{babel}
\usepackage[T1]{fontenc}
\usepackage{amsmath}
\usepackage{amsfonts}
\usepackage{amssymb}
\usepackage{graphicx}
\usepackage{listings}
\usepackage{xcolor}
\usepackage{hyperref}

\title{RTOS Dokumentation: Danielou OS (DOS)}
\author{Danielou Mounsande}
\date{\today}

\definecolor{mGreen}{rgb}{0,0.6,0}
\definecolor{mGray}{rgb}{0.5,0.5,0.5}
\definecolor{mPurple}{rgb}{0.58,0,0.82}
\lstset{
    commentstyle=\color{mGreen},
    keywordstyle=\color{magenta},
    numberstyle=\tiny\color{mGray},
    stringstyle=\color{mPurple},
    basicstyle=\footnotesize\ttfamily,
    breaklines=true,
    captionpos=b,
    keepspaces=true,
    numbers=left,
    showspaces=false,
    showstringspaces=false,
    showtabs=false,
    tabsize=2,
    language=C,
    extendedchars=true,
    literate={ü}{{\"u}}1 {ä}{{\"a}}1 {ö}{{\"o}}1 {ß}{{\ss}}1
}

\begin{document}

\maketitle
\tableofcontents
\newpage

\section{Einleitung}
Diese Dokumentation dient als technisches Referenzwerk für die Entwicklung des eigenen Echtzeitbetriebssystems (DOS) auf dem STM32L475. Sie beschreibt die Architektur, die Funktionszyklen und dokumentiert die Lösungen für kritische Probleme.

\section{Technische Definitionen (Glossar)}

\subsection{TCB (Task Control Block)}
Der TCB ist der \textbf{Personalausweis} eines Tasks. In DOS wird er durch die Struktur \texttt{TCB\_sctTCB\_t} definiert.
\begin{itemize}
    \item \textbf{u32TaskSP:} Der aktuelle Stack Pointer des Tasks. Er liegt an Offset 0 für schnellen Assembler-Zugriff.
    \item \textbf{eTaskState:} Zustand des Tasks (Ready, Running, Blocked).
    \item \textbf{u32TicksToWait:} Verbleibende System-Ticks, die eine Task im Blocked-Zustand verharren muss.
    \item \textbf{au32TaskStack:} Reservierter Speicherbereich für den Task-Stack.
\end{itemize}

\subsection{Stack Frame (Stapelrahmen)}
\begin{itemize}
    \item \textbf{Hardware Stack Frame:} Wird vom Cortex-M4 automatisch gesichert (xPSR, PC, LR, R12, R3, R2, R1, R0).
    \item \textbf{Software Stack Frame:} Wird vom Kernel manuell gesichert (R4 bis R11).
\end{itemize}

\subsection{System-Register \& Hardware-Steuerung}
\begin{itemize}
    \item \textbf{CPACR (Coprocessor Access Control Register):} Steuert den Zugriff auf Coprozessoren. Wir nutzen es, um die FPU (Floating Point Unit) zu deaktivieren. Dies vereinfacht den Kontextwechsel, da wir sonst 32 zusätzliche Register (S0-S31) sichern müssten.
    \item \textbf{SHPR3 (System Handler Priority Register 3):} Hier setzen wir die Priorität von \textbf{PendSV} auf den niedrigsten Wert (0xFF). Das sorgt dafür, dass ein Task-Wechsel niemals kritische Hardware-Interrupts blockiert.
    \item \textbf{PendSV:} Ausnahme mit niedrigster Priorität für den eigentlichen Kontextwechsel.
    \item \textbf{SysTick:} Taktgeber für die präemptive Zeitsteuerung (Heartbeat).
    \item \textbf{ICSR (Interrupt Control State Register):} Ein Steuerregister des Cortex-M, über das wir den PendSV-Interrupt per Software auslösen können.
\end{itemize}

\subsection{AAPCS \& Register R0}
\begin{itemize}
    \item \textbf{AAPCS:} Der \textit{Procedure Call Standard for ARM} definiert die Regeln für Funktionsaufrufe.
    \item \textbf{Register R0 (Die Schnittstelle):} R0 ist besonders, da es laut Standard zwei Rollen hat:
    \begin{enumerate}
        \item \textbf{Input:} Das erste Argument einer C-Funktion wird in R0 übergeben.
        \item \textbf{Output:} Der Rückgabewert einer Funktion wird in R0 zurückgegeben.
    \end{enumerate}
    In unserem Kernel nutzen wir R0, um den Stack-Pointer zwischen Assembler und C zu übergeben.
\end{itemize}

\subsection{Alignment (Ausrichtung)}
Gemäß AAPCS muss der Stack Pointer bei Funktionsaufrufen auf \textbf{8 Byte (64-Bit)} ausgerichtet sein. Dies ist notwendig für Befehle wie \texttt{LDRD/STRD} und die FPU. Ohne dieses Alignment kann es zu \texttt{UsageFaults} kommen.

\section{Funktionszyklen}

\subsection{Task-Erstellungszyklus}
\begin{enumerate}
    \item \textbf{Allokation:} TCB aus dem globalen Array entnehmen.
    \item \textbf{Initialisierung (Dummy Stack Frame):}
    \begin{itemize}
        \item PC auf Task-Funktion setzen + Thumb-Bit (LSB=1).
        \item xPSR Bit 24 auf 1 setzen.
        \item LR auf \texttt{Kernel\_TaskReturnHook} setzen.
    \end{itemize}
    \item \textbf{Bereitstellung:} Zustand auf \texttt{TaskState\_Ready} setzen.
\end{enumerate}

\subsection{Kontextwechsel-Zyklus (Scheduling)}
1. \textbf{Trigger:} SysTick setzt das PendSV-Bit im ICSR. \\
2. \textbf{Assembler (PendSV):} Sichert R4-R11, ruft \texttt{Kernel\_ContextSwitch} auf. \\
3. \textbf{C-Logik:} Sichert alten SP im TCB, wählt neuen Task (Round Robin), gibt neuen SP zurück. \\
4. \textbf{Wiederherstellung:} Assembler lädt R4-R11 vom neuen Stack, aktualisiert MSP, \texttt{bx lr} stellt Hardware-Frame wieder her.

\section{Forensische Analyse: Die großen Probleme}

\subsection{Der INVPC (0x00040000) Absturz}
\textbf{Problem:} HardFault mit Invalid PC load.\\
\textbf{Ursache:} Das LSB (Least Significant Bit) des Task-PC war 0. ARM Cortex-M unterstützt nur den \textbf{Thumb-Modus}. Wenn das LSB 0 ist, versucht der Prozessor in den ARM-Modus zu wechseln, was zum Absturz führt.\\
\textbf{Lösung:} Erzwingen von \texttt{address | 0x01} bei der Task-Erstellung sowie ein Cast nach \texttt{(uint32_t)}, um Bit-Operationen auf Adressen zu ermöglichen.

\subsection{Null-Pointer Korruption \& Pipeline-Optimierung}
\textbf{Problem:} Der erste Kontextwechsel benötigt eine Sonderbehandlung, da noch kein Task läuft. Früher wurde dies über einen \texttt{cbz}-Check im Assembler gelöst.\\
\textbf{Lösung:} Verschiebung der Initialisierungs-Logik in die C-Ebene (\texttt{Kernel\_ContextSwitch}). Der Assembler-Teil bleibt linear und performant ("Autobahn-Prinzip"), während C entscheidet, ob ein SP gespeichert werden muss.

\section{FAQ: Häufig gestellte Fragen}

\subsection{Warum nutzen wir nur den MSP?}
Zur Vereinfachung. In professionellen Systemen trennt man MSP (Kernel) und PSP (User), um den Kernel vor Stack-Overflow der Tasks zu schützen.

\subsection{Warum R4-R11 manuell sichern?}
Das Hardware-Auto-Stacking sichert nur R0-R3, R12, LR, PC und xPSR. Ein RTOS muss den kompletten Zustand (inkl. R4-R11) bewahren, damit lokale Variablen der Tasks nicht korrumpiert werden.

\subsection{Warum \texttt{\_\_attribute__((naked))}?}
Dieser Attribut verhindert, dass der Compiler automatisch Prolog- und Epilog-Code (Stack-Manipulation) für den \texttt{PendSV\_Handler} generiert. Ohne dies würde der Compiler den Stack verändern, bevor wir den exakten CPU-Zustand sichern können.

\subsection{Was bewirkt \texttt{nop} im Idle-Task?}
\texttt{nop} (No Operation) verhindert, dass der Compiler eine leere Endlosschleife wegoptimiert. Es hält die CPU in einem kontrollierten Leerlauf.

\subsection{Warum Dekrementierung (\texttt{--sp}) beim Stack-Aufbau?}
Der Cortex-M nutzt einen \textbf{Full Descending Stack}. Wir müssen den Pointer erst verringern, um in den gültigen Speicherbereich des Arrays zu gelangen, bevor wir einen Wert schreiben (sonst droht ein Buffer Overflow).

\section{Kernel Objekte \& Zeitsteuerung}

\subsection{Blocking-State \& Kernel\_TaskDelay}
Um CPU-Ressourcen zu schonen, wurde ein blockierendes Warten implementiert.
\begin{enumerate}
    \item \textbf{Aufruf:} Eine Task ruft \texttt{Kernel\_TaskDelay(ticks)} auf.
    \item \textbf{Zustandswechsel:} Der Kernel setzt \texttt{eTaskState} auf \texttt{TaskState\_Blocked} und speichert die Zeit in \texttt{u32TicksToWait}.
    \item \textbf{Yield:} Ein sofortiger Kontextwechsel wird erzwungen, damit andere Tasks laufen können.
    \item \textbf{Update:} Im \texttt{SysTick\_Handler} wird jede Millisekunde \texttt{Kernel\_UpdateTimers()} aufgerufen.
    \item \textbf{Wakeup:} Wenn \texttt{u32TicksToWait} Null erreicht, wird die Task zurück auf \texttt{TaskState\_Ready} gesetzt und vom Scheduler wieder berücksichtigt.
\end{enumerate}

\subsection{Semaphore: Synchronisation und Ressourcenschutz}
Semaphoren ermöglichen die Koordination von Tasks. In DOS wurde eine effiziente Blockier-Logik implementiert.

\subsubsection{Der "Wait"-Zyklus}
Wenn eine Task \texttt{Kernel\_SemaphoreWait} aufruft:
\begin{enumerate}
    \item \textbf{Kritischer Abschnitt:} Interrupts werden deaktiviert (\texttt{cpsid i}), um atomaren Zugriff auf den Zähler zu garantieren.
    \item \textbf{Prüfung:} Wenn \texttt{u32Count > 0}, wird der Zähler dekrementiert und die Task läuft sofort weiter.
    \item \textbf{Blockierung:} Wenn \texttt{u32Count == 0}:
    \begin{itemize}
        \item Die Adresse des Semaphors wird im TCB-Feld \texttt{pWaitingObject} gespeichert.
        \item Der Zustand wechselt auf \texttt{TaskState\_Blocked}.
        \item Ein Kontextwechsel wird angefordert.
    \end{itemize}
\end{enumerate}

\subsubsection{Der "Give"-Zyklus}
Wenn eine Task \texttt{Kernel\_SemaphoreGive} aufruft:
\begin{enumerate}
    \item \textbf{Inkrementierung:} Der Zähler wird erhöht (bis zum definierten \texttt{u32MaxCount}).
    \item \textbf{Gezielte Suche:} Der Kernel durchläuft alle TCBs und sucht nach einer Task, deren \texttt{pWaitingObject} exakt auf die Adresse dieses Semaphors zeigt.
    \item \textbf{Wakeup:} Die gefundene Task wird auf \texttt{TaskState\_Ready} gesetzt und \texttt{pWaitingObject} wird gelöscht (\texttt{NULL}).
    \item \textbf{Dispatch:} Ein Kontextwechsel wird ausgelöst, damit die aufgeweckte Task ggf. sofort übernehmen kann.
\end{enumerate}

\subsection{Mutexe: Gegenseitiger Ausschluss mit Ownership}
Ein Mutex ist ein spezielles binäres Semaphor mit dem Konzept des Besitzes (\textbf{Ownership}). Nur die Task, die den Mutex gesperrt hat, darf ihn wieder freigeben.

\subsubsection{Ownership-Logik}
\begin{itemize}
    \item Beim \texttt{Lock} wird im Mutex-Objekt ein Zeiger \texttt{pOwner} auf den aktuellen TCB gespeichert.
    \item Beim \texttt{Unlock} prüft der Kernel: \texttt{if (pMutex->pOwner == Global\_PointerToCurrentlyRunningTCB)}.
    \item Dies verhindert, dass Tasks fälschlicherweise Ressourcen freigeben, die sie nicht kontrollieren.
\end{itemize}

\subsection{Priority Inheritance (Prioritätsvererbung)}
Um das Problem der unbeschränkten \textbf{Prioritätsinversion} zu lösen, wurde ein vollständiges Prioritätsvererbungs-Protokoll (PIP) im Kernel implementiert.

\subsubsection{Vererbungs-Mechanismus (Lock)}
Wenn eine hochpriore Task $H$ einen Mutex sperren will, der von einer niedrigprioren Task $L$ gehalten wird:
\begin{enumerate}
    \item \textbf{Detektion:} Der Kernel stellt fest, dass der Mutex gesperrt ist (\texttt{u8IsLocked == 1}) and ermittelt den Besitzer (\texttt{pOwner}).
    \item \textbf{Vergleich:} Falls die aktuelle Priorität von $H$ höher ist (kleinere Zahl) als die des Besitzers $L$, wird die Priorität von $L$ temporär angehoben:
    \begin{align*}
        L.\text{u8CurrentPriority} = H.\text{u8CurrentPriority}
    \end{align*}
    \item \textbf{Transitive Vererbung (Kettenbildung):} Falls die Task $L$ ihrerseits blockiert ist und auf einen anderen Mutex wartet, der von einer Task $M$ gehalten wird, breitet sich die Prioritätsvererbung rekursiv über die Kette aus. Die Task $M$ erbt ebenfalls die Priorität von $H$.
\end{enumerate}

\subsubsection{Rückvererbungs-Mechanismus (Unlock)}
Wenn die Task $L$ den Mutex wieder freigibt:
\begin{enumerate}
    \item \textbf{Prioritätsberechnung:} Die Task $L$ darf nicht blind auf ihre Basis-Priorität (\texttt{u8BasePriority}) zurückfallen, falls sie noch \textbf{andere} Mutexe hält, auf die andere hochpriore Tasks warten.
    \item \textbf{Scannen:} Der Kernel durchsucht alle TCBs im System. Für jede blockierte Task, die auf einen Mutex wartet, der sich im Besitz von $L$ befindet, wird deren Priorität ausgewertet.
    \item \textbf{Wiederherstellung:} Die neue Priorität von $L$ wird auf das Maximum (höchste Priorität / kleinste Zahl) aus ihrer Basis-Priorität und den Prioritäten aller auf sie wartenden Tasks gesetzt.
    \item \textbf{Bereitstellung:} Der am höchsten priorisierte wartende Thread auf dem freigegebenen Mutex wird aufgeweckt.
\end{enumerate}

\section{Message Queues: Inter-Task-Kommunikation}
Zur sicheren Übertragung von Daten und Synchronisation zwischen Tasks wurden Message Queues implementiert.

\subsection{Struktur und Pufferung}
Eine Message Queue (\texttt{Kernel\_Queue\_t}) basiert auf einem Ringpuffer mit fester maximaler Kapazität (\texttt{KERNEL\_QUEUE\_MAX\_SIZE} = 16 Elemente).
Die Verwaltung erfolgt über folgende Felder:
\begin{itemize}
    \item \textbf{au32Buffer:} Ein Array aus 32-Bit-Ganzzahlen, das Rohdaten oder Zeiger (Pointer) auf komplexe Strukturen aufnehmen kann.
    \item \textbf{u8Head / u8Tail:} Schreib- und Lese-Zeiger im Ringpuffer.
    \item \textbf{u8Count / u8Size:} Aktuelle Anzahl der Nachrichten und maximale Kapazität.
\end{itemize}

\subsection{Send- und Receive-Betriebsmodi}
Die Kernel-API bietet für Lese- und Schreiboperationen sowohl blockierende als auch nicht-blockierende Aufrufe:

\subsubsection{Schreiben (Send)}
\begin{itemize}
    \item \textbf{Puffer frei:} Die Nachricht wird in den Puffer geschrieben, \texttt{u8Head} inkrementiert (modulo \texttt{u8Size}), and die am höchsten priorisierte Task, die auf Daten wartet, wird aufgeweckt.
    \item \textbf{Puffer voll (Nicht-Blockierend):} Die Funktion bricht sofort ab und liefert den Fehlercode \texttt{KernelError\_QueueFull}.
    \item \textbf{Puffer voll (Blockierend):} Die schreibende Task wird in den Zustand \texttt{TaskState\_Blocked} versetzt and das \texttt{pWaitingObject} auf die Queue gesetzt. Sie wird erst reaktiviert, wenn ein Consumer Daten liest.
\end{itemize}

\subsubsection{Lesen (Receive)}
\begin{itemize}
    \item \textbf{Nachricht vorhanden:} Die Nachricht wird aus dem Puffer gelesen, \texttt{u8Tail} inkrementiert, and die am höchsten priorisierte Task, die auf freien Platz wartet, wird aufgeweckt.
    \item \textbf{Puffer leer (Nicht-Blockierend):} Die Funktion bricht sofort ab und liefert den Fehlercode \texttt{KernelError\_QueueEmpty}.
    \item \textbf{Puffer leer (Blockierend):} Die lesende Task blockiert und wartet, bis ein Producer Daten sendet.
\end{itemize}

\section{Referenzen}
\begin{itemize}
    \item PM0214: STM32 Cortex-M4 Programming Manual.
    \item AAPCS Standard.
    \item Kursunterlagen Echtzeitsysteme (THM SS26).
\end{itemize}

\end{document}
